
\chapter*{Laboratorio computacional: de las ecuaciones al reporte}
\phantomsection\addcontentsline{toc}{chapter}{Laboratorio computacional: R y Python}\label{lab:comp}
\begin{idea}[Pregunta de trabajo]
\emph{¿Una tabla de transiciones estimada con lenguaje técnico asigna mayor
probabilidad a bitácoras técnicas reservadas que una tabla estimada con prosa general?}
Responder esta pregunta construida en pequeño ayuda a comprender qué medir,
pero no demuestra nada sobre un tokenizador moderno ni sobre la población real
de documentos industriales.
\end{idea}
\section*{L1. Inspeccionar el corpus antes de ajustar el modelo}
El archivo \nolinkurl{experimentos/corpus_didactico.csv} contiene ocho documentos
generales de entrenamiento, ocho técnicos de entrenamiento y seis técnicos
de prueba; todas las frases fueron \textbf{redactadas para esta clase}.
\textbf{No se descargaron informes de empresas ni se entrenó un LLM.}
Las dos condiciones usan el mismo vocabulario de caracteres establecido
\emph{solo con entrenamiento} y el mismo valor $\alpha=1$; los datos
de prueba nunca intervienen en el ajuste ni en la selección del modelo.
Se usan minúsculas y se quitan espacios de los extremos; el archivo se almacena
en UTF-8 con caracteres precompuestos, sin normalización adicional.

\section*{L2. Traducir ecuaciones a un algoritmo comprobable}
Para cada documento $z=z_1\cdots z_n$, se añade un estado inicial
\texttt{BOS} y se predice también un símbolo \texttt{EOS}. La probabilidad
predictiva de bigramas es
\[
\widehat P(v\mid u)=\frac{C(u,v)+\alpha}
{\sum_w C(u,w)+\alpha|V_{\mathrm{salida}}|},\qquad \alpha=1.
\]
La pérdida en bits para el documento completo es
\[
B(z)=-\sum_{t=1}^{n+1}\log_2 \widehat P(z_t\mid z_{t-1}),
\qquad \operatorname{BPC}(z)=\frac{B(z)}{n}.
\]
Aquí $z_0=\mathrm{BOS}$ y $z_{n+1}=\mathrm{EOS}$. El denominador
cuenta los puntos de código del texto original (\emph{no} EOS).
\textbf{Observación:} el alfabeto es el mismo y los modelos difieren en
los datos de entrenamiento, no en el algoritmo de tokenización.

\section*{L3. Ejecutar en dos lenguajes y localizar los resultados}
Desde el directorio que contiene \texttt{s01\_integrada.tex}:
\begin{verbatim}
python codigo/representacion.py
python codigo/s01.py
python codigo/experimento_s01.py
# En un equipo donde este instalado R:
Rscript codigo/representacion.R
Rscript codigo/s01.R
Rscript codigo/experimento_s01.R
\end{verbatim}
Los programas nuevos exportan los resultados a
\nolinkurl{experimentos/resultados_python.csv} y
\nolinkurl{experimentos/resultados_R.csv}. Python exporta también
\nolinkurl{experimentos/resumen_python.json}. Si falta \texttt{Rscript},
la ejecución R se deja \textbf{PENDIENTE}; no se informa una comparación
entre lenguajes como completada.

\section*{L4. Interpretar un resultado ejecutado, no extrapolarlo}
\begin{ejemplo}[Salida verificada por Python en este corpus sintético]
Se obtuvieron seis diferencias pareadas
$\Delta_i=\operatorname{BPC}_{\rm tecnico}(z_i)-
\operatorname{BPC}_{\rm general}(z_i)$; su promedio fue
$\bar\Delta\approx-0.54025$ bits por carácter.
Por enumeración exhaustiva de las $6^6=46656$ remuestras bootstrap
pareadas con reemplazo, el intervalo percentil nominal del 95\% fue
$[-0.66771,-0.42038]$.
El intervalo \emph{describe el comportamiento de un procedimiento de
remuestreo sobre seis frases construidas}, no la incertidumbre científica
de una población de reportes: no hay muestreo representativo, comparación
con modelos neuronales ni validación externa.
\end{ejemplo}
\section*{L5. De algoritmo a evidencia: registros obligatorios}
El reporte contiene la versión del programa, datos y licencia,
partición por documento, protocolo BOS/EOS, alfabeto, unidades,
valores por documento, método del intervalo, limitaciones y
una conclusión que no exceda lo medido. Cada cambio del corpus,
semillas, modelado o tokenización requiere regenerar todas las tablas.
El \hyperref[lab:ic10]{microexperimento IC10} evalua, por separado,
la segmentacion aprendida por dos tokenizadores BPE.
\ifdefined\labcompstandalone
La consigna del taller exige un reporte corto; el seminario de
investigación explica que falta
\else
La \hyperref[mod:taller]{consigna del taller} exige un reporte corto;
el \hyperref[mod:invest]{seminario de investigación} explica qué falta
\fi
para convertir una demostración en un estudio candidato a publicación.
